\begin{abstract}
Let (A) and (B) be involutions and suppose their ordered product
(H=AB) has order (2n), where (n\geq3). We prove that the generated
group is the dihedral group of order (4n), that the reverse product is
(H^{-1}=BA), and that the encounter commutator is

\[
c=[A,B]=H^2.
\]

The element (c) has order (n) and generates the commutator subgroup.
The unique order-two element of the lifted cyclic subgroup is the binary
return receipt
\[
\tau=H^n.
\]
Their relationship is controlled exactly by the parity of the visible cycle
length. If (n) is even, then (	au=c^{n/2}), so the commutator tower reaches
the binary return receipt. If (n) is odd, then
(	au\notin\langle c\rangle),
\[
\langle H\rangle=\langle c\rangle\times\langle\tau\rangle,
\qquad
H=c^{(n+1)/2}\tau,
\]
and the commutator and return receipt are independent cyclic factors. In the
natural Klein-four encounter quotient, (	au) maps to the identity exactly
for even (n); for odd (n) it maps to the visible product class. Thus cycle
doubling is universal, while concealment of the binary receipt by the
encounter quotient is parity-sensitive. The order-eight, order-sixteen
receipt tower is recovered at (n=4). No physical interpretation is claimed.
\end{abstract}

